\subsection{Listas y conjuntos}
\label{alg:2-7}

La lista es la estructura por defecto en Python y el conjunto es su complemento para
preguntar "¿está?" rápido. Casi todos los errores de tiempo con listas vienen de
usar una operación \bigO{n} dentro de un ciclo sin notarlo.

\subsubsection{Crear listas}

\begin{lstlisting}[language=Python]
n, m = 3, 4
a = [0] * n                        # [0, 0, 0]
g = [[0] * m for _ in range(n)]    # matriz n x m con filas independientes
c = list(range(1, n + 1))          # [1, 2, 3]
b = [x * x for x in c if x > 1]    # [4, 9]  comprension con filtro
\end{lstlisting}

\graybold{Cuidado:} \texttt{[[0] * m] * n} NO crea una matriz: repite la misma fila
$n$ veces, y cambiar una casilla cambia toda la columna.

\vspace{0.3cm}
\subsubsection{Ordenar}

\begin{lstlisting}[language=Python]
a = [3, 1, 2]
a.sort()                           # ordena la misma lista, devuelve None
b = sorted(a, reverse=True)        # lista nueva: [3, 2, 1]
p = [(1, 5), (2, 3), (3, 5)]
p.sort(key=lambda t: (-t[1], t[0]))  # por t[1] descendente; empate: t[0] asc.
# p == [(1, 5), (3, 5), (2, 3)]
\end{lstlisting}

Las tuplas se comparan campo por campo, así que una llave de tupla ordena por
varios criterios a la vez; negar un campo numérico lo invierte. El error típico es
\texttt{a = a.sort()}, que deja \texttt{a} en \texttt{None}.

\vspace{0.3cm}
\subsubsection{Costo de las operaciones}

\begin{center}
\begin{tabular}{|l|l|}
\hline
Operación & Costo \\ \hline
\texttt{a[i]}, \texttt{len(a)}, \texttt{a.append(x)}, \texttt{a.pop()} & \bigO{1} \\ \hline
\texttt{a.pop(0)}, \texttt{a.insert(0, x)} & \bigO{n} \\ \hline
\texttt{x in a}, \texttt{a.index(x)}, \texttt{a.count(x)} & \bigO{n} \\ \hline
\texttt{a[l:r]} (crea una copia) & \bigO{r - l} \\ \hline
\texttt{x in s} con \texttt{s} conjunto & \bigO{1} \\ \hline
\end{tabular}
\end{center}

Para sacar por el frente se usa \texttt{deque}; para preguntar pertenencia muchas
veces, un \texttt{set}.

\vspace{0.3cm}
\subsubsection{min, max y sum}

\begin{lstlisting}[language=Python]
pal = ["aa", "b", "cccc"]
max(pal, key=len)                  # 'cccc'
min([(1, 9), (2, 4)], key=lambda t: t[1])   # (2, 4)
max([], default=0)                 # 0; sin default da ValueError
sum(x for x in range(10) if x % 2) # 25
\end{lstlisting}

\vspace{0.3cm}
\subsubsection{Imprimir listas}

\begin{lstlisting}[language=Python]
a = [1, 2, 3]
print(*a)                          # 1 2 3
print(*a, sep="\n")                # uno por linea
print(*a, sep="")                  # 123
\end{lstlisting}

Es cómodo para salidas cortas; con miles de líneas conviene
\texttt{sys.stdout.write} con \texttt{join} (sección 1).

\vspace{0.3cm}
\subsubsection{Conjuntos}

\begin{lstlisting}[language=Python]
s = {1, 2, 3}
s.add(4)
s.discard(9)                       # no falla si no esta (remove si falla)
a, b = {1, 2, 3}, {2, 3, 4}
a | b, a & b, a - b, a ^ b         # union, interseccion, diferencia, simetrica
len(set([5, 5, 7]))                # 2: cantidad de distintos
visto = frozenset(a)               # inmutable: sirve como llave de un dict
\end{lstlisting}

Un \texttt{set} no tiene orden: para imprimirlo ordenado, \texttt{sorted(s)}.

\vspace{0.3cm}
\subsubsection{Copiar}

\begin{lstlisting}[language=Python]
import copy
a = [1, 2, 3]
b = a[:]                           # copia (igual que a.copy() o list(a))
g = [[1, 2], [3, 4]]
g2 = [fila[:] for fila in g]       # copia de matriz: rapida
g3 = copy.deepcopy(g)              # copia de cualquier anidamiento: mas lenta
\end{lstlisting}

\texttt{b = a} no copia nada: son dos nombres para la misma lista. Y \texttt{a[:]}
sobre una matriz copia solo la lista exterior; las filas siguen compartidas.
